% GENERATED FILE. Edit canonical lessons, then run npm run generate:books.
% canonical-source-hash: fnv1a64:e6bc0ea7c928acbe
% canonical-lessons: ML-C74-aan, ML-C74-purpose, ML-C74-kazhiyum, ML-R74-aan-recall

\chapter{One Ending, Two Sentences}
\label{ch:ml-aan}
\hlchaptermodality{74}

\begin{chapteropening}
\textbf{By the end of this chapter:} \emph{I can say what I am going somewhere in order to do, and say what I am able to do.}

From cold: the -uka to -aan swap, the purpose sentence it builds in front of another verb, and the ability sentence it builds in front of kazhiyum — with the person in the dative, the way knowing already works.
\end{chapteropening}

\section[-āṉ]{-\ml{ആൻ} — one ending swapped for another}

\label{lesson:ML-C74-aan}



\begin{quote}
Say the verb for \emph{to read}. Then say two more verbs from that chapter, and listen to how all of them end.

That shared ending is about to be replaced.
\end{quote}

\subsection*{What to know first}
\begin{quote}
\textbf{-\ml{ആൻ}} — \emph{-āṉ} — \textbf{the \textquotedblleft{}to do it\textquotedblright{} form}
\end{quote}

Every verb you have met ends in \textbf{-\ml{ഉക}} when it is sitting in a dictionary. Take that ending off and put \textbf{-\ml{ആൻ}} on instead:

\noindent
\begin{tabularx}{\linewidth}{@{}>{\raggedright\arraybackslash}X>{\raggedright\arraybackslash}X@{}}
\toprule
\textbf{the dictionary form} & \textbf{the -\ml{ആൻ} form} \\
\midrule
\textbf{\ml{വായിക്കുക}} \emph{vāyikkuka} & \textbf{\ml{വായിക്കാൻ}} \emph{vāyikkāṉ} \\
\textbf{\ml{പോകുക}} \emph{pōkuka} & \textbf{\ml{പോകാൻ}} \emph{pōkāṉ} \\
\textbf{\ml{കാണുക}} \emph{kāṇuka} & \textbf{\ml{കാണാൻ}} \emph{kāṇāṉ} \\
\bottomrule
\end{tabularx}

\textbf{Nothing else in the word changes.} The piece in front of the ending — the stem — stays exactly as it was, so a verb you can say you can already swap.

This is the ending Malayalam builds two separate things on, and the next two lessons are those two things. Getting the swap into your hand now makes both of them free.

\subsection*{Your turn}
Say these aloud:

\begin{itemize}
\raggedright
  \item \emph{vāyikkuka} … \emph{vāyikkāṉ}
  \item \emph{pōkuka} … \emph{pōkāṉ}
  \item \emph{kāṇuka} … \emph{kāṇāṉ}
  \item which part of each word stayed still
\end{itemize}

\begin{tcolorbox}[breakable,colback=teal!4,colframe=teal!35!black,title={Before you move on}]
Which ending comes off? (\textbf{-\ml{ഉക}}.) Which goes on? (\textbf{-\ml{ആൻ}}.) And what happens to the rest of the word? (\textbf{Nothing} — the stem does not move.)
\end{tcolorbox}
\section[vāyikkāṉ pōkunnu]{\ml{വായിക്കാൻ} \ml{പോകുന്നു} — going in order to do something}

\label{lesson:ML-C74-purpose}



\begin{quote}
Swap the ending on three verbs: \emph{to read}, \emph{to go}, \emph{to see}.

Now put one of those in front of another verb and you have said why you are doing it.
\end{quote}

\subsection*{What to know first}
\begin{quote}
\textbf{\ml{ഞാൻ} \ml{വായിക്കാൻ} \ml{പോകുന്നു}.}
\end{quote}

\begin{quote}
\emph{ñān vāyikkāṉ pōkunnu} — \textquotedblleft{}I am going \textbf{in order to read}.\textquotedblright{}
\end{quote}

The \textbf{-\ml{ആൻ}} form comes \textbf{first}, and the verb that actually happens comes \textbf{last}. Read it in that order and it says: \emph{I — to-read — go.}

\textbf{English puts the purpose after the verb} and often needs a word for it — \emph{I am going to read}, \emph{I came to see you}. Malayalam puts the purpose in front and needs no extra word at all: the ending is the whole of it.

\noindent
\begin{tabularx}{\linewidth}{@{}>{\raggedright\arraybackslash}X>{\raggedright\arraybackslash}X@{}}
\toprule
\textbf{} & \textbf{} \\
\midrule
\textbf{\ml{കാണാൻ} \ml{വന്നു}} \emph{kāṇāṉ vannu} & came in order to see \\
\textbf{\ml{പറയാൻ} \ml{വന്നു}} \emph{paṟayāṉ vannu} & came in order to say \\
\bottomrule
\end{tabularx}

\begin{grammarlens}[title={Grammar Lens: the order is the meaning}]
Malayalam has put the important verb last all along — the quotative chapter said the same thing about a reported sentence, where the saying came at the end. \textbf{This is that habit again.} The reason goes in front, the event goes at the back, and the reader hears what you did last.

That order is worth expecting rather than learning per construction: when a sentence in this language starts with something that is not a verb you can finish, keep going — the verb you need is at the end.
\end{grammarlens}

\subsection*{Your turn}
Say these aloud:

\begin{itemize}
\raggedright
  \item \emph{ñān vāyikkāṉ pōkunnu}
  \item \emph{kāṇāṉ vannu}
  \item which of the two verbs in each sentence is the thing that happened
  \item where English would put the purpose that Malayalam puts first
  \item the dictionary form of the going-verb, and then its -\ml{ആൻ} form
\end{itemize}

\begin{tcolorbox}[breakable,colback=teal!4,colframe=teal!35!black,title={Before you move on}]
Which comes first, the purpose or the event? (\textbf{The purpose}.) How do you say \textquotedblleft{}I am going in order to read\textquotedblright{}? (\emph{\textbf{Ñān vāyikkāṉ pōkunnu}}.) And how many extra words does Malayalam need for \emph{in order to}? (\textbf{None} — the ending does it.)
\end{tcolorbox}
\section[kaḻiyuṁ]{\ml{കഴിയും} (kaḻiyuṁ) — \textquotedblleft{}can\textquotedblright{}, built the way knowing is built}

\label{lesson:ML-C74-kazhiyum}



\begin{quote}
Say \emph{enikku malayāḷaṁ aṟiyāṁ}. Say which word in it means \emph{I}, and what shape that word is in.

Today's sentence is built the same way.
\end{quote}

\subsection*{What to know first}
\begin{quote}
\textbf{\ml{എനിക്ക്} \ml{വായിക്കാൻ} \ml{കഴിയും}.}
\end{quote}

\begin{quote}
\emph{enikku vāyikkāṉ kaḻiyuṁ} — \textquotedblleft{}I can read.\textquotedblright{}
\end{quote}

Three pieces, and \textbf{two of them you already had}: \emph{enikku}, the dative \emph{to me} that the knowing sentence uses, and \emph{vāyikkāṉ}, the \textbf{-\ml{ആൻ}} form from two pages ago. The new word is \textbf{\ml{കഴിയും}}, and what it says is that the thing is possible.

\textbf{Ability arrives at you in Malayalam. It is not something you do.} The knowing chapter said this in as many words — knowing arrives, so \emph{I} leaves the subject slot and turns into \emph{to me}. Can behaves identically:

\noindent
\begin{tabularx}{\linewidth}{@{}>{\raggedright\arraybackslash}X>{\raggedright\arraybackslash}X@{}}
\toprule
\textbf{} & \textbf{} \\
\midrule
\textbf{\ml{എനിക്ക്} \ml{മലയാളം} \ml{അറിയാം}} & Malayalam is known to me \\
\textbf{\ml{എനിക്ക്} \ml{വായിക്കാൻ} \ml{കഴിയും}} & reading is possible to me \\
\bottomrule
\end{tabularx}

English makes both of these an action with \emph{I} at the front — \emph{I know}, \emph{I can}. Malayalam puts the person in the dative and lets the ability be the subject.

\begin{grammarlens}[title={Grammar Lens: one ending, two jobs}]
The \textbf{-\ml{ആൻ}} form has now done both of the things this chapter promised:

\noindent
\begin{tabularx}{\linewidth}{@{}>{\raggedright\arraybackslash}X>{\raggedright\arraybackslash}X@{}}
\toprule
\textbf{} & \textbf{} \\
\midrule
\textbf{\ml{വായിക്കാൻ} \ml{പോകുന്നു}} & going \textbf{in order to} read \\
\textbf{\ml{വായിക്കാൻ} \ml{കഴിയും}} & \textbf{able to} read \\
\bottomrule
\end{tabularx}

Same ending, same swap, two different sentences — and the word after it decides which one you are saying.
\end{grammarlens}

\subsection*{Your turn}
Say these aloud:

\begin{itemize}
\raggedright
  \item \emph{enikku vāyikkāṉ kaḻiyuṁ}
  \item the knowing sentence and the can sentence one after the other
  \item what the two of them share at the front
  \item \emph{pōkāṉ kaḻiyuṁ} — being able to go
\end{itemize}

\begin{tcolorbox}[breakable,colback=teal!4,colframe=teal!35!black,title={Before you move on}]
How do you say \emph{I can read}? (\emph{\textbf{Enikku vāyikkāṉ kaḻiyuṁ}}.) What shape is the word for \emph{I} in? (\textbf{The dative} — \emph{to me}.) And which sentence you already knew is built the same way? (\textbf{The one about knowing Malayalam}.)
\end{tcolorbox}
\section[āṉ ōrmma]{(one ending, two sentences) — from cold}

\label{lesson:ML-R74-aan-recall}



\begin{quote}
Close the lessons before this one. Take \emph{to read}, \emph{to go} and \emph{to see} out of the dictionary and swap their endings.
\end{quote}

\begin{grammarlens}[title={Grammar Lens: the swap, and what it buys}]
\textbf{One swap: -\ml{ഉക} off, -\ml{ആൻ} on, stem untouched.}

\noindent
\begin{tabularx}{\linewidth}{@{}>{\raggedright\arraybackslash}X>{\raggedright\arraybackslash}X@{}}
\toprule
\textbf{what follows the -\ml{ആൻ} form} & \textbf{what the sentence says} \\
\midrule
a verb that happened & \textbf{in order to} — \emph{vāyikkāṉ pōkunnu} \\
\textbf{\ml{കഴിയും}} & \textbf{can} — \emph{enikku vāyikkāṉ kaḻiyuṁ} \\
\bottomrule
\end{tabularx}

Two sentences from one ending, and the word after it decides which. \textbf{Neither needed a new verb}: every verb in this chapter came out of the reading and going chapters, and the only genuinely new word is \emph{kaḻiyuṁ}.

The ability sentence also puts the person in the \textbf{dative}, the way the knowing sentence does. That is the second time this language has said the same thing: some things are done by you, and some arrive at you — and it marks which is which rather than leaving you to guess.
\end{grammarlens}

\subsection*{Your turn}
Say these aloud:

\begin{itemize}
\raggedright
  \item the swap on three verbs
  \item \emph{ñān vāyikkāṉ pōkunnu}
  \item \emph{enikku vāyikkāṉ kaḻiyuṁ}
  \item which of those two has the person in the dative, and why
\end{itemize}

\begin{tcolorbox}[breakable,colback=teal!4,colframe=teal!35!black,title={Before you move on}]
Which ending comes off and which goes on? (\textbf{-\ml{ഉക} off, -\ml{ആൻ} on}.) What does the -\ml{ആൻ} form mean in front of another verb? (\textbf{In order to}.) And what does it mean in front of \emph{kaḻiyuṁ}? (\textbf{Can}.)
\end{tcolorbox}
